
\subsubsection{No long path in $\Dcontr$}\label{sec:biclique:nolongpath}

In this section, we prove Lemma~\ref{lem:digging} in the case when the length of a longest path in $\Dcontr$ is $p \leq 2t+2$. Recall that $\Dcontr$ contains $(A,B)$-matching and $\Dcontr^{\star}$ contains $(A,B)$-induced-biclique. Also, the edges of $A \cup B$ are minimal in $\Dcontr$. We prove Lemma~\ref{lem:digging} by induction on the length of a longest $(A,B)$-path in $\Dcontr$.
For this purpose we essentially restate Lemma~\ref{lem:digging} in a form that is ``induction-friendly''. Observe, as a base case, that when the length of a longest $(A,B)$-path is $1$, then since $\Dcontr^{\star}$ contains $(A,B)$-induced-biclique, we conclude that $\Dcontr$ contains $(A,B)$-induced-biclique and the edges of the $(A,B)$-induced-biclique are minimal. In this case, we conclude that we get the second output of Lemma~\ref{lem:digging}.

\begin{lemma}\label{lem:biclique:nolongpath}
Let $A,B \subseteq V(\Dcontr)$ be ordered sets such that the length
of a longest $(A,B)$-path in $\Dcontr$ is $p$,
$|A|= |B| \geq f(t,p)$, $\Dcontr^{\star}$ contains an $(A,B)$-induced-biclique and the matching edges of $(A,B)$ are minimal in $\Dcontr$, then either
\begin{enumerate}
\item there exist ordered sets $A',B' \subseteq V(\Dcontr)$ such that the length
of a longest $(A',B')$-path in $\Dcontr$ is at most $p$,
$|A'|= |B'| \geq f(t,p-1)$, $\Dcontr^{\star}$ contains an $(A',B')$-induced-biclique and the matching edges of $(A',B')$ are minimal in $\Dcontr$, or
\item $\Dcontr$ contains an ordered $t$-\toughpair, or
\item $\Dcontr^{\star}$ contains a $t$-induced-biclique whose edges are minimal in $\Dcontr$.
\end{enumerate}
\end{lemma}

\begin{proof}
Let $L$ be the set of vertices of $\Dcontr$ that contains the first internal vertex on every longest $(A,B)$-path in $\Dcontr$.

\begin{claim}\label{claim:biclique6}
There is no $(L,L)$-path in $\Dcontr$.
\end{claim}

\begin{claimproof}
For the sake of contradiction, suppose there exist $x,y \in L$ such that there this an $(x,y)$-path in $\Dcontr$. Note that $x\neq y$, as otherwise, $\Dcontr$ would contain a cycle.
From the definition of $L$, there exists $a_i,a_j \in L$ ($i$ not necessarily distinct from $j$) such that $x$ is the first internal vertex on some $(a_i,B)$-path, say $P_x$, of length $p$ and $y$ is the first internal vertex on some $(a_j,B)$-path, say $P_y$, of length $p$. 
Let $P_{xy}$ be some $(x,y)$-path in $\Dcontr$. We first claim that the set of internal vertices of $P_{x,y}$ is disjoint from $P_y$. Suppose not. Let $z$ be some internal vertex of $P_{xy}$ that is also a vertex of $P_y$. Since $z \in V(P_y)$, there exists a $(y,z)$-path in $\Dcontr$. Also since $z \in V(P_{x,y})$, there exists a $(z,y)$-path in $\Dcontr$. This implies a cycle in $\Dcontr$, which is a contradiction. Thus, we conclude that the internal vertices of $P_{xy}$ are disjoint from that of $P_y$. Then consider the $(x,B)$-path obtained by appending $P_{xy}$ and $P_y$. Note that the length of $P_y$ is $p-1$ and the length of $P_{xy}$ is at least $1$. This $(x,B)$-path, together with the edge $(a_i,x)$, gives an $(a_i,B)$-path of length strictly greater than $p$ in $\Dcontr$, which contradicts that the length of a longest $(A,B)$-path in $\Dcontr$ is $p$.
\end{claimproof}

\begin{claim}\label{claim:biclique7}
For each $(a,x) \in E(\Dcontr)$ such that $a \in A$ and $x \in L$, $(a,x)$ is a minimal edge of $\Dcontr$. 
\end{claim}

\begin{claimproof}
For the sake of contradiction, say there exists an $(a,x)$-path in $\Dcontr$ of length at least $2$. Let this path be $P_{ax}$. Since $x \in L$, there exists an $(A,B)$-path of length $p$ whose first internal vertex is $x$. Let this path be $P_{AB}$. Then the internal vertices of $P_{ax}$ are disjoint from the internal vertices of $P_{AB}$, as otherwise there would be a cycle in $\Dcontr$. By appending the path $P_{sx}$ (which is of length at least $2$) with the subpath of $P_{ABs}$ starting from $x$ (which is of length $p-1$), we get an $(A,B)$-path in $\Dcontr$ of length at least $p+1$, which is a contradiction.
\end{claimproof}

\begin{claim}\label{claim:biclique8}
Let $a \in A$ and $x \in L$, such that $(a,x) \not \in E(\Dcontr)$. Then $(a,x) \not \in E(\Dcontr^{\star})$.
\end{claim}

\begin{claimproof}
For the sake of contradiction, suppose there exists an $(a,x)$-path, say $P_{a,x}$ in $\Dcontr$ of length at least two. Since $x \in L$, there exists $a' \in A$ ($a'$ not necessarily distinct from $a$) such that $x$ is the first internal vertex of a $(a',T)$-path, say $P_{a'B}$, of length $p$. First observe that the internal vertices of $P_{a,x}$ are disjoint from that of $P_{a'B}$, as otherwise there would be a cycle in $\Dcontr$. The path $P_{a,x}$ appended with the $(x,B)$-subpath of $P_{a'B}$ gives an $(A,B)$-path in $\Dcontr$ of length strictly greater than $p$, which is a contradiction.
\end{claimproof}

Consider the bipartite graph $D_{\bip}=\Dcontr[A \cup L]$. Note, from Claim~\ref{claim:biclique7}, that each edge of $D_{\bip}$ is a minimal edge in $\Dcontr$.
We now distinguish into two cases based on the size of the matching in $D_{\bip}$.\\

\noindent{\bf Case 1: The size of a maximum matching in $D_{\bip}$ is at least $\mathcal{R}'(2t)$.} In this case, from Ramsey's Theorem, either there exists an induced matching of size $2t$ in $D_{\bip}$ or a $2t$-induced-biclique. From Claim~\ref{claim:biclique7}, each edge of $D_{\bip}$ is a minimal edge of $D_{contr}$. 

In the first case, we get sets $A^* \subseteq A$ and $B^* \subseteq L$ such that there is an $(A^*,B^*)$-induced-matching in $D_{\bip}$. Since $A$ is an independent set in $\Dcontr^{\star}$, so is $A^*$. Since $B^* \subseteq L$, from Claim~\ref{claim:biclique7}, $B^*$ is an independent set in $\Dcontr^{\star}$. Thus, $(A^*,B^*)$ is an induced-matching in $\Dcontr$. Further, from Claim~\ref{claim:biclique8}, $(A^*,B^*)$ is in fact an induced matching in $\Dcontr^{\star}$. Observe that if a digraph has an induced-matching of size $2t$ in its transitive closure, then it has a ordered $t$-\toughpair. Thus, in this case we get the second outcome of the lemma.

In the second case, when $H_{\bip}$ contains a $2t$-induced-biclique, then this is also a $2t$-induced-biclique in $D_{\contr}$ where all edges of the biclique are minimal edges of $\Dcontr$. Thus, in this case we get the third outcome of the lemma.\\

\noindent{\bf Case 2: The size of a maximum matching in $D_{\bip}$ is at most $\mathcal{R}'(2t)$.} In this case, 
we will find ordered sets $A',B' \subseteq V(\Dcontr)$ satisfies the properties stated in the first outcome of the lemma.

Let $Z$ be a minimum vertex cover of $D_{\bip}$ of size at most $\mathcal{R}'(2t)$. Let $Z_A = Z \cap A$ and 
$Z_L=Z \cap L$. Let $A^*=A \setminus Z_S$. Since $Z_A \cup Z_L$ is a vertex cover of $D_{\bip}$, for any  $a \in A^*$, $N_L(a) \subseteq Z_L$.
For each subset $L' \subseteq Z_L$, let $A_{L'} \subseteq A^*$ such that for each vertex $a \in A_{L'}$, $N_{Z_L}(a)=L'$ (and hence $N_L(a)=L'$).

Fix $L' \subseteq Z_L$ such that $|S_{L'}|$ is maximized. 
Recall that the sets $A =(a_1,\ldots,a_{f(t,p)})$ and $T=(t_1, \ldots, t_{f(t,p)})$ are ordered sets.
Set $A'=A_{L'}$ and $B'$ be the corresponding vertices of $B$, 
that is $b_i \in B'$ if and only if $a_i \in A'$.
From the choice of $L'$, $|A'|=|B'| \geq (|A|-\mathcal{R}'(2t))/2^{\mathcal{R}'(2t)} = f(t,p-1)$.
We will now show that the length of a longest $(A',B')$-path in $\Dcontr$ is at most $p-1$. 

\begin{claim}\label{claim:biclique5}
Let $a_i \in A'$ and $b_j \in B'$ such that $i \neq j$. Let $P=(a_i,x_1,\ldots,x_q,b_j)$ be a longest $(A',B')$-path in $\Dcontr$. Then $x_1 \not \in L$.
\end{claim}
\begin{claimproof}
Since $b_j \in B'$, $a_j \in A'$ (from construction of $B'$).
First observe that since $b_j \in B'$, then $a_j \in A'$.
For the sake of contradiction, suppose $x_1 \in L$. Since $(a_i,x) \in E(\Dcontr)$ and all the vertices of $A'$ have the same neighbourhood in $Z_L$ and $a_j \in A'$, we conclude that $(a_j,x) \in E(\Dcontr)$. Thus, there is an $(a_j,b_j)$-path in $\Dcontr$ of length strictly greater than one, which contradicts that $(a_j,b_j)$ is a minimal edge of $\Dcontr$.
\end{claimproof}

From Claim~\ref{claim:biclique5}, the first internal vertex of every longest $(A',B')$-path is not contained in $L$. Now suppose that there exists an $(A',B')$-path of length $p$. Then its first internal vertex should be in $L$ by the definition of $L$. This is a contradiction.  
\end{proof}

Lemma~\ref{lem:biclique:nolongpath} proves Lemma~\ref{lem:digging} when the length of the longest path $p \leq 2t+2$.